Thresholding is one of the most straightforward and widely used approaches for detecting blink activity from EEG or EOG signals ~\citep{tran2021detection,chang2016detection}.
The simplicity of thresholding makes it computationally efficient and easy to implement, particularly in real-time fatigue-detection systems~\citep{tran2021detection,jefri2022eye}.
%In addition, thresholding is often used as an initial detection step in more complex pipelines, where the detected candidates are further refined using morphological rules or machine-learning classifiers ~\citep{jiang2023eyeblink, ghosh2023automatic}.
Existing approaches range from single- or few-channel detectors that apply amplitude thresholds with rule-based decision logic to hybrid pipelines that combine thresholding as an initial detection step with additional candidate-refinement mechanisms ~\citep{nolan2010faster, chang2016detection,jiang2023eyeblink, ghosh2023automatic,tran2021detection,zhang2023method,wang2022multidimensional,wang2025sliding,kleifges2017blinker,cao2021unsupervised,agarwal2019blink,valderrama2018automatic,guttmann2019new}.
%[8, 9, 19, 24, 25].
Despite these developments, selecting an appropriate threshold remains challenging because blink amplitude varies across participants, electrode locations, recording devices, signal quality, and experimental conditions ~\citep{guttmann2019new,wang2022multidimensional}.
Blink amplitude can vary across participants, electrode locations, recording devices, signal quality, and experimental conditions~\citep{guttmann2019new,wang2022multidimensional}.
A threshold that is too low may produce false detections due to noise or other artifacts, whereas a threshold that is too high may miss genuine blink events~\citep{tran2021detection,jefri2022eye}.
This backdrop explains why simple one-value global thresholds can be attractive, yet brittle.
Therefore, adaptive and robust threshold selection remains a critical consideration for reliable blink detection, particularly in fatigue-driving applications~\citep{ko2020eyeblink,shahbakhti2023fusion}.

